\chapter{Projeto final: seu primeiro livro}
\label{chap:projeto-final}

Chegou a hora de transformar as técnicas anteriores num projeto autoral. Use o
modelo da pasta \arquivo{examples/} como ponto de partida ou copie este próprio
projeto e remova o conteúdo didático.

\section{Etapa 1: definir o livro}

Antes do LaTeX, escreva uma frase para cada item:

\begin{itemize}
  \item tema e transformação prometida ao leitor;
  \item público e conhecimento prévio necessário;
  \item entre cinco e dez capítulos, cada um com um resultado observável;
  \item exemplos, imagens ou fontes de pesquisa necessários.
\end{itemize}

Um sumário provisório é uma hipótese, não um contrato. Ele pode mudar quando a
escrita revelar uma ordem mais clara.

\section{Etapa 2: montar a estrutura}

Crie \arquivo{main.tex}, \arquivo{preamble.tex}, \arquivo{bibliography.bib} e as
pastas \arquivo{chapters/} e \arquivo{images/}. Faça um arquivo por capítulo e
adicione todos ao principal com \comando{include}. Compile enquanto os capítulos
ainda contêm apenas títulos. Assim, problemas de infraestrutura são resolvidos
antes de se misturarem a centenas de páginas.

\section{Etapa 3: escolher a identidade visual}

Defina uma fonte de leitura, uma fonte de títulos e uma fonte monoespaçada se
houver código. Escolha uma cor principal e uma cor de destaque. Teste uma página
com título, parágrafos, lista, negrito e itálico. Priorize legibilidade; efeitos
visuais devem apoiar a hierarquia, não competir com o conteúdo.

\section{Etapa 4: escrever e revisar}

Produza primeiro uma versão completa, mesmo imperfeita. Depois faça revisões
separadas:

\begin{enumerate}
  \item \textbf{estrutura}: ordem e objetivo de cada capítulo;
  \item \textbf{conteúdo}: precisão, exemplos e lacunas;
  \item \textbf{linguagem}: clareza, repetição e terminologia;
  \item \textbf{visual}: quebras, figuras, tabelas e consistência;
  \item \textbf{técnica}: log limpo, links, citações e índice.
\end{enumerate}

\section{Lista de verificação}

Seu primeiro livro está pronto para uma versão de leitura quando:

\begin{itemize}
  \item compila do zero com um único comando;
  \item não possui referências ou citações indefinidas;
  \item explica a quem se destina e como deve ser usado;
  \item mantém títulos, exemplos e termos consistentes;
  \item possui fontes e imagens com licenças adequadas;
  \item foi lido no PDF, não apenas no código-fonte;
  \item foi testado por pelo menos uma pessoa do público pretendido.
\end{itemize}

\begin{resultado}
  Ao concluir essas etapas, você terá um livro modular, compilável e preparado
  para crescer. Mais importante: saberá localizar cada decisão editorial no
  código e poderá alterá-la sem refazer o documento inteiro.
\end{resultado}

\section*{Próximos passos}

Explore a documentação oficial do projeto LaTeX \citep{latex-project} à medida
que surgirem necessidades concretas. Aprender pacotes sob demanda é mais
eficiente do que tentar memorizar o ecossistema. Seu primeiro livro não precisa
usar todos os recursos; precisa servir bem ao leitor.
